% !TeX root = ../动手学强化学习笔记.tex
% 由主文档通过 \input 载入；请编译主文档。
\setcounter{chapter}{20} % 保留原书章号。
\chapter{分层强化学习与时间抽象}
\label{ch:rl-hierarchy}

分层强化学习（hierarchical reinforcement learning）把跨多个时间步的行为组织成可调用技能。高层决定做什么，低层决定如何执行；收益来自合适的抽象和复用，而不是层次越多就必然越好。本章承接多步回报与规划，以给定选项和技能库为主要分析对象，建立半马尔可夫备份并讨论技能组合的最优性范围。技能学习部分说明子目标、内部奖励与终止条件的关系，自动技能发现算法留待进一步学习。

\section{选项与技能表示}

选项（option）由启动集合 $\mathcal I_o$、内部策略 $\pi_o$ 和终止概率 $\beta_o(s)$ 组成。在允许的状态启动选项后，低层按 $\pi_o$ 逐步执行动作，每一步后依终止规则决定是否交回高层。高层策略 $\nu(o\mid s)$ 只在决策时刻选择可启动选项。

“走到门口”可以是选项，单步“向右移动”可以看作立即在一步后终止的特殊选项。为方便本章分析，假设每个启动的选项至少执行一步且几乎必然在有限时间终止；真正任务终止也会结束当前选项。

\section{半马尔可夫备份}

选项执行时长 $K$ 随起点、环境和内部策略变化，因此高层决策构成半马尔可夫决策过程（semi-Markov decision process，SMDP）。固定高层策略下，选项价值满足
\[
q_\nu(s,o)=\mathbb E\!\left[\sum_{k=0}^{K-1}\gamma^kR_{k+1}
+\gamma^K\sum_{o'}\nu(o'\mid S_K)q_\nu(S_K,o')\,\middle|\,S_0=s,o\right].
\]
若后继是真实终止状态，最后价值项取零。高层最优备份将后继的策略平均换成可启动选项上的最大值。折扣按原始时间步计算，必须保留随机 $\gamma^K$，不能把任意长选项都只乘一次 $\gamma$。

\begin{example}[持续时间改变价值]
设 $\gamma=1/2$，两个选项结束后均到达价值为 $4$ 的状态。第一选项持续一步，奖励为 $1$，目标为 $1+(1/2)4=3$；第二选项持续两步，奖励为 $0,3$，目标为 $0+(1/2)3+(1/4)4=2.5$。虽然后者未折扣累计奖励更高，高层仍可能偏好较快完成的前者。
\end{example}

\section{技能学习、子目标与最优性范围}

技能可人工指定，也可围绕子目标训练。例如在多房间网格中，为各门口训练到达策略，高层再选择经过哪些房间。技能评价应关注可启动范围、到达概率、持续时间和失败后状态，而不只是子任务平均奖励。

只允许从固定技能库组合行为，会限制可执行的原始动作序列。因此“高层已经最优”通常只表示在该技能库内最优，不等于原始 MDP 全局最优。若所有原始动作都作为一步选项保留，表达限制会减弱，但学习和规划仍可能复杂。技能内部奖励也可能与外部任务冲突，终止条件应与高层需求匹配。

\section{习题与实验}
\begin{enumerate}
\item 推导固定持续时间 $K=2$ 时的选项备份，并说明随机持续时间为何必须保留期望内部的 $\gamma^K$。
\item 构造一个技能库，使最优原始动作序列无法表示；再加入一个技能或原始动作修复它。
\item 在多房间环境比较原始动作控制与固定技能高层控制，同时统计成功率、原始交互步数和高层决策次数，不把不同时间尺度的“步数”直接比较。
\end{enumerate}





